\section{If the rate can depend on the source exponent}
The guarantee derived above uses a policy whose rates work throughout \(0<\varepsilon<1/52\). The source itself fixes \(\varepsilon\) before selecting a witness. Could a policy that reads this parameter improve on the answer? The order of optimization matters in general. Here the two surviving supports make it possible to solve the fixed-parameter version for the literal prefactor test. This supplies a check on whether an additional uniformity convention is doing substantial work.

Put
\[
 n_r=6r+1,\quad \beta=3r/n_r,\quad
 h=2+8\varepsilon,\quad \ell=\frac{1}{n_r}\log\frac{h+2}{h+3}.
\]
The strict logarithmic cell test at this exponent is \(\log10-\mu+\beta\log\nu>\ell\). Define \(C_\varepsilon\) as the large root of the leaf equation and \(B_\varepsilon\) as the small root of the sparse-parent equation
\begin{align*}
 \log10-C_\varepsilon+\beta\log(C_\varepsilon/2)&=\ell,\\
 \log10-\frac{B_\varepsilon+5C_\varepsilon}{32}
                           +\beta\log(B_\varepsilon/32)&=\ell .
\end{align*}
On the local ranges of interest the first left-hand side decreases in \(C_\varepsilon\), the second increases in \(B_\varepsilon\). The worst-label supremum for the stationary, exact-count, edge-realizing class certified by this test at that \(\varepsilon\) is
\[
 F_\varepsilon=
 \frac{116C_\varepsilon}{2112+254C_\varepsilon+32B_\varepsilon}.       \tag{X}
\]
The formula describes the continuum of exponents, including the jumps in \(r=\lceil32/\varepsilon+64\rceil+1\). Interpolation of a small numerical grid would miss the integer changes.

To justify the upper bound simultaneously for this range, start from
\(\log(4/5)/10381<\ell<0\) and \(5190/10381\le\beta<1/2\). Exact logarithmic comparisons enclose
\[
 2.39208<C_\varepsilon<2.39213,\qquad
 0.7058<B_\varepsilon<0.7064 .
\]
For instance the leaf expression at \(2.39208\) with the smallest \(\beta\) exceeds zero; at \(2.39213\) even its largest endpoint falls below the lower bound on \(\ell\). For the sparse equation use the lower cap, lower singleton, and smaller \(\beta\) to bound its largest possible value at the lower singleton boundary. At the upper singleton boundary do the reverse. The smaller exponent maximizes the expression here because \(\nu<1\). These monotonicities avoid an uncontrolled interchange of the two endpoint exponents.

At the root \(\mu=1\). If only one preferred leaf is available, \(\nu\le 2.39213/64\); even the largest possible log margin at that value falls below \(\log(4/5)/10381\). Thus the root still requires at least two preferred leaves. The near-threshold census was rerun at the enlarged rational cap \(239213/100000\); the same two supports survive. Its root-hole residual is strictly negative. For each of the two masks the half's preferred singleton must meet a lower rate \(b\ge h_\varepsilon(s)\), where \(s\) denotes its five companions' rate sum. Implicit differentiation of the sparse equation gives
\[
 h_\varepsilon'(s)=
   \frac{h_\varepsilon(s)}{32\beta-h_\varepsilon(s)}<1/15
       \qquad (0\le s\le5C_\varepsilon).
\]
The last bound follows from the positive small root, its monotonic increase, the displayed rectangle, and the lower bound on \(\beta\). Hence the surplus comparison from the two-mask calculation is valid with \(C,B\) replaced by \(C_\varepsilon,B_\varepsilon\). The first mask's eight extra denominator units can only reduce its upper; the second gives (X).

For attainment in the limiting sense, on the frozen exceptional support put 26 rates just below \(C_\varepsilon\), put the singleton just above its corresponding lower root, and set the adjustment by the count. At a fixed \(\varepsilon\) the two defining equalities themselves are excluded by the strict test. Bounds uniform on the displayed rectangle verify that all the other 21 hard forms have log slack above \(0.14\) even under the standardized test. Moving into the strict region and then approaching the corner proves the lower supremum. The ordinary lower guarantee exceeds the largest possible (X) on the rectangle. In particular no new support policy is needed to attain the fixed-exponent answer to arbitrary precision.

\begin{center}\small\begin{tabular}{@{}lrr@{}}\toprule
Chosen \(\varepsilon\)& \(r\)&Worst-label supremum under (W), percent\\\midrule
\(64/3329\)&1730&10.1191110389\\
\(1/64\)&2113&10.1190959370\\
\(1/256\)&8257&10.1190436792\\
\(1/1024\)&32833&10.1190298549\\\bottomrule
\end{tabular}\end{center}
Each numerical entry was separately enclosed by rational bisection of both equations. These examples are not needed to establish (X). They show the size of the advantage at specific admissible exponents and test a point at the minimum source exponent.

Finally let \(\varepsilon\downarrow0\). Then \(\ell\to0,\ \beta\to1/2\), and the uniformly enclosed roots converge to \(C_*,B_*\). A rule allowed to read \(\varepsilon\) still cannot promise a larger fraction \emph{simultaneously} at every exponent, since its putative floor would exceed the optimal value at sufficiently small \(\varepsilon\). The common-rate construction already supplies the limiting guarantee. Thus the distinction matters for a particular preassigned exponent, but it disappears for the single robust number requested across the full family. This last statement concerns test (W); the raw scalar selector can admit stationary profiles by other arguments.
